% Toque MCIC sobre IEEEtran: solo color y diagramas. No cambia márgenes,
% fuentes, columnas ni tamaños de la plantilla IEEE.
% tectonic compila con XeTeX: Times New Roman (la de IEEE y la secundaria del manual UD)
% se lee de las fuentes del manual de marca; si no está, TeX Gyre Termes (clon de Times).
\usepackage{amsmath,amssymb}
\usepackage{newtxmath}
\usepackage{fontspec}
\providecommand{\UDmarca}{/Users/juferoga/repos/ud/maestria/mcic/.UniversidadDistrital-manual-marca}
\IfFileExists{\UDmarca/fuentes/times.ttf}{%
  \setmainfont{times.ttf}[Path=\UDmarca/fuentes/, BoldFont=timesbd.ttf,
    ItalicFont=timesi.ttf, BoldItalicFont=timesbi.ttf]}{%
  \setmainfont{texgyretermes}[Extension=.otf, UprightFont=*-regular, BoldFont=*-bold,
    ItalicFont=*-italic, BoldItalicFont=*-bolditalic]}
\setmonofont{texgyrecursor}[Extension=.otf, UprightFont=*-regular, BoldFont=*-bold, Scale=0.95]
\usepackage[spanish,es-noshorthands,es-tabla,es-nolists]{babel}
\usepackage{graphicx}
\usepackage{booktabs,tabularx,array,multirow}
\usepackage[table]{xcolor}
\usepackage{tikz}
\usetikzlibrary{positioning,arrows.meta,calc,fit}
\usepackage[hidelinks]{hyperref}
\usepackage{cite}
\usepackage{url}

% Paleta del manual de marca UD (design.md)
\definecolor{udazul}{HTML}{1291C0}
\definecolor{udazulosc}{HTML}{0C6F96}
\definecolor{uddorado}{HTML}{B49739}
\definecolor{udverde}{HTML}{00A859}
\definecolor{udplata}{HTML}{98999A}
\definecolor{udsable}{HTML}{373435}
\definecolor{udrojo}{HTML}{ED3237}

\hypersetup{colorlinks=true,linkcolor=udazulosc,citecolor=udazulosc,urlcolor=udazulosc}

% Bloques de la ecuación de búsqueda: el mismo color en ecuación, PRISMA y matriz.
\newcommand{\bloqueA}[1]{\textcolor{udazulosc}{#1}}
\newcommand{\bloqueB}[1]{\textcolor{uddorado!85!black}{#1}}
\newcommand{\bloqueC}[1]{\textcolor{udverde!80!black}{#1}}
\newcommand{\etiqueta}[2]{\tikz[baseline=(t.base)]\node[fill=#1,text=white,inner sep=1.6pt,
  rounded corners=1pt,font=\bfseries\scriptsize](t){#2};}

% Cabecera de tablas: filete azul y fila sombreada clara
\newcommand{\cabecera}{\rowcolor{udazul!12}}
\arrayrulecolor{udazulosc}

% Columnas de la matriz
\newcolumntype{L}[1]{>{\raggedright\arraybackslash}p{#1}}
\newcolumntype{Y}{>{\raggedright\arraybackslash}X}

\makeatletter
\providecommand{\filas}[1]{\@@input #1 }
\makeatother

% Marca de enfoque en la matriz (A..E)
\newcommand{\enf}[1]{\textcolor{udazulosc}{\bfseries #1}}

% Cajas del PRISMA
\tikzset{
  caja/.style={draw=udazulosc,line width=0.5pt,rounded corners=1.5pt,align=center,
    font=\scriptsize,inner sep=2.5pt,text width=#1},
  caja/.default=3.6cm,
  excl/.style={draw=udplata,line width=0.5pt,rounded corners=1.5pt,align=left,
    font=\scriptsize,inner sep=2.5pt,fill=udplata!8,text width=#1},
  excl/.default=2.75cm,
  fase/.style={fill=udazul,text=white,font=\bfseries\tiny,minimum width=0.42cm,
    rounded corners=1pt,inner sep=1pt},
  flecha/.style={-{Latex[length=1.6mm]},udsable,line width=0.5pt},
}
